% lp-manual-style.tex: the Landing Party field manual's page and components (LuaLaTeX, A5 booklet).
% Read into the preamble by lp-manual.lua (which defines \lpextdir, \lptextwidthmm and the factions first).
% Every macro here is a look; every value a page shows comes from the document or its data file.
\makeatletter

% ---------------------------------------------------------------- page
% A5, two-sided. Inner margin a little wider than the outer: the fold eats some, the thumb tabs sit outside.
% The text block is \lptextwidthmm wide (the tree layout in lp-manual.lua uses the same number).
\RequirePackage{geometry}
\geometry{paperwidth=148mm, paperheight=210mm, textwidth=\lptextwidthmm mm, inner=16mm,
  top=18mm, bottom=20mm, headheight=10pt, headsep=5mm, footskip=10mm, twoside}
\RequirePackage{eso-pic}
\RequirePackage{tikz}
\usetikzlibrary{calc,arrows.meta,positioning,fadings}
\RequirePackage{fancyhdr}
\RequirePackage{titlesec}
\RequirePackage{tcolorbox}
\tcbuselibrary{skins}
\RequirePackage{lettrine}
\RequirePackage{enumitem}
\RequirePackage{etoolbox}
\RequirePackage{microtype}
\RequirePackage{colortbl}

% ---------------------------------------------------------------- type
% Body: Alegreya (TeX Live, OFL). A literary serif built for long text, sturdy enough on screen at A5 sizes and
% with real small caps and old-style figures: the narrator's voice. Titillium Web (Factorio's own face, OFL,
% vendored in fonts/) for the machine's voice: headings, stat panels, labels, captions, trees.
\RequirePackage{fontspec}
\setmainfont{Alegreya}[Extension=.otf, UprightFont=*-Regular, BoldFont=*-Bold, ItalicFont=*-Italic,
  BoldItalicFont=*-BoldItalic, FontFace={m}{sc}{AlegreyaSC-Regular}, FontFace={b}{sc}{AlegreyaSC-Bold},
  Numbers={OldStyle,Proportional}]
\setsansfont{TitilliumWeb}[Path=\lpextdir/fonts/, Extension=.ttf, UprightFont=*-Regular, BoldFont=*-Bold,
  ItalicFont=*-Regular, ItalicFeatures={FakeSlant=0.16}, BoldItalicFont=*-Bold, BoldItalicFeatures={FakeSlant=0.16},
  FontFace={sb}{n}{*-SemiBold}]
\newfontfamily\lpsemibold{TitilliumWeb-SemiBold}[Path=\lpextdir/fonts/, Extension=.ttf]
\newfontfamily\lplining{Alegreya}[Extension=.otf, UprightFont=*-Regular, ItalicFont=*-Italic, Numbers=Lining]
\renewcommand\normalsize{\@setfontsize\normalsize{9.6}{13}\abovedisplayskip 8pt plus 2pt\belowdisplayskip 8pt plus 2pt}
\renewcommand\small{\@setfontsize\small{8.8}{11.6}}
\renewcommand\footnotesize{\@setfontsize\footnotesize{8}{10.4}}
\renewcommand\scriptsize{\@setfontsize\scriptsize{7}{8.8}}
\normalsize
\newcommand\lplabel[1]{{\lpsemibold\addfontfeature{LetterSpace=9}\MakeUppercase{#1}}}
\newcommand\lpkicker[1]{{\lpsemibold\fontsize{6.6}{8}\selectfont\addfontfeature{LetterSpace=14}\MakeUppercase{#1}}}
\newcommand\HeaderFont{\sffamily\bfseries}% the clip extension's placeholder asks for it
\input{\lpextdir/lp-manual-emblems.tex}
\setlength{\parindent}{0pt}
\setlength{\parskip}{4pt plus 1.5pt minus 1pt}
\setlength{\emergencystretch}{2.5em}
\widowpenalty=10000 \clubpenalty=10000 \raggedbottom

% ---------------------------------------------------------------- colour
% Chrome is neutral and prints on any paper; colour is the factions' (defined from the data by the filter).
\definecolor{LPInk}{HTML}{1C1E22}
\definecolor{LPNight}{HTML}{14161A}
\definecolor{LPMuted}{HTML}{6B6E74}
\definecolor{LPRule}{HTML}{C9C5BC}
\definecolor{LPPaper}{HTML}{F5F0E6}
\definecolor{LPPaperDark}{HTML}{E7DFD0}
\color{LPInk}
% The kinan house names that other extensions (clip) draw with.
\colorlet{AccentPrimary}{LPMuted}\colorlet{AccentLight}{LPPaper}\colorlet{TextSecondary}{LPMuted}

% \lpdefinefaction{id}{HEX}{on: white|ink}{tab index, 0 for none}{display name}
\newcommand\lpdefinefaction[5]{%
  \definecolor{lp@#1}{HTML}{#2}%
  \colorlet{lp@#1@dark}{lp@#1!58!black}%
  \colorlet{lp@#1@light}{lp@#1!9!white}%
  \colorlet{lp@#1@mid}{lp@#1!28!white}%
  \ifstrequal{#3}{white}{\colorlet{lp@#1@on}{white}}{\colorlet{lp@#1@on}{LPInk}}%
  \expandafter\def\csname lp@#1@tab\endcsname{#4}%
  \expandafter\def\csname lp@#1@name\endcsname{#5}}
\lpdefinefaction{neutral}{8A8E94}{white}{0}{}
% \lpusefaction{id}: the current faction (chapter, card, box). Fac, FacDark, FacLight, FacMid, FacOn.
\newcommand\lpusefaction[1]{%
  \ifcsname lp@#1@tab\endcsname\def\lp@fac{#1}\else\def\lp@fac{neutral}\fi
  \colorlet{Fac}{lp@\lp@fac}\colorlet{FacDark}{lp@\lp@fac @dark}\colorlet{FacLight}{lp@\lp@fac @light}%
  \colorlet{FacMid}{lp@\lp@fac @mid}\colorlet{FacOn}{lp@\lp@fac @on}%
  \edef\lp@factab{\csname lp@\lp@fac @tab\endcsname}\edef\lp@facname{\csname lp@\lp@fac @name\endcsname}}
\newcommand\lpchapterfaction{neutral}
\lpusefaction{neutral}

% ---------------------------------------------------------------- page furniture
% Thumb tabs: a faction's chapters carry a tab on the fore-edge at the faction's own height, so the folded
% booklet can be opened at a faction by its edge. Lore pages are printed on a paper tint.
\newif\iflplore
\newif\iflpnotab
\newcommand\lp@tab{%
  \begingroup\lpusefaction{\lpchapterfaction}%
  \ifnum\lp@factab>0\relax\iflpnotab\else
    \begin{tikzpicture}[x=1mm,y=1mm,remember picture=false]
      \useasboundingbox (0,0) rectangle (148,210);
      \pgfmathsetmacro\lp@tabtop{190-(\lp@factab-1)*27}
      \ifodd\value{page}
        \fill[Fac] (144.5,\lp@tabtop) rectangle (151,\lp@tabtop-24);
        \node[rotate=-90, font=\lpsemibold\fontsize{5.2}{6}\selectfont\addfontfeature{LetterSpace=12}, text=FacOn]
          at (146.4,\lp@tabtop-12) {\MakeUppercase{\lp@facname}};
      \else
        \fill[Fac] (-3,\lp@tabtop) rectangle (3.5,\lp@tabtop-24);
        \node[rotate=90, font=\lpsemibold\fontsize{5.2}{6}\selectfont\addfontfeature{LetterSpace=12}, text=FacOn]
          at (1.6,\lp@tabtop-12) {\MakeUppercase{\lp@facname}};
      \fi
    \end{tikzpicture}%
  \fi\fi\endgroup}
% Period dress, after the manuals of the 1990s (StarCraft's race-textured frames and the faint drawing behind its unit
% pages, C&C's edge tabs): every page on a paper ground (generated, media/ornaments, make_paper.sh); a chapter's pages
% carry its faction's emblem as a watermark (the filigrane) and corner pieces in its colour; a blank page carries the
% book's own seal. The watermark stays under 8% of the faction's dark tone: body text over it keeps its contrast, and a
% greyscale print shows it as a faint line.
\newif\iflpdecor   % from the first chapter to the back cover
\newif\iflpblank   % this page is a blank (set by the clears and the padding)
\newif\iflp@paper \newif\iflp@lorepaper
\newsavebox\lp@paperbox \newsavebox\lp@lorebox
% \lpsetpaper{body ground}{lore ground}: image paths, either empty for none (the filter gives them). Each is set in a
% saved box once, so the PDF holds one copy however many pages print it.
\newcommand\lpsetpaper[2]{%
  \ifstrempty{#1}{}{\global\lp@papertrue\global\setbox\lp@paperbox\hbox{\includegraphics[width=\paperwidth,height=\paperheight]{#1}}}%
  \ifstrempty{#2}{}{\global\lp@lorepapertrue\global\setbox\lp@lorebox\hbox{\includegraphics[width=\paperwidth,height=\paperheight]{#2}}}}
\newcommand\lp@ground{%
  \iflplore\iflp@lorepaper\copy\lp@lorebox\else{\color{LPPaper}\rule{\paperwidth}{\paperheight}}\fi
  \else\iflp@paper\copy\lp@paperbox\fi\fi}
\newcommand\lpwatermarkopacity{0.075}
\newcommand\lp@watermark{%
  \begingroup\lpusefaction{\lpchapterfaction}%
  \begin{tikzpicture}[x=1mm,y=1mm]
    \useasboundingbox (0,0) rectangle (148,210);
    \clip (0,0) rectangle (148,210);
    \ifodd\value{page}\def\lp@wx{110}\else\def\lp@wx{38}\fi
    \begin{scope}[shift={(\lp@wx,74)}, transparency group, opacity=\lpwatermarkopacity, color=FacDark]
      \lpemblem{\lp@fac}{58}{1.25}
    \end{scope}
  \end{tikzpicture}\endgroup}
% Corner pieces, 7 mm in from the trim: a bracket in the faction's colour, its finish the faction's own.
\newcommand\lp@corners{%
  \begingroup\lpusefaction{\lpchapterfaction}%
  \begin{tikzpicture}[x=1mm,y=1mm]
    \useasboundingbox (0,0) rectangle (148,210);
    \foreach \cx/\cy/\sx/\sy in {7/7/1/1, 141/7/-1/1, 7/203/1/-1, 141/203/-1/-1}{%
      \begin{scope}[shift={(\cx,\cy)}, xscale=\sx, yscale=\sy]
        \ifcsname lp@corner@\lp@fac\endcsname\csname lp@corner@\lp@fac\endcsname\else\lp@corner@neutral\fi
      \end{scope}}
  \end{tikzpicture}\endgroup}
\tikzset{lpcornerA/.style={draw=Fac, line width=1.15pt, line cap=butt, line join=miter},
  lpcornerB/.style={draw=FacDark, line width=0.3pt}}
\def\lp@corner@neutral{\draw[lpcornerA] (0,12) -- (0,0) -- (12,0); \draw[lpcornerB] (1.7,8.5) -- (1.7,1.7) -- (8.5,1.7);
  \fill[Fac] (0,0) +(0,1.3) -- +(1.3,0) -- +(0,-1.3) -- +(-1.3,0) -- cycle;}
\expandafter\def\csname lp@corner@kessler-hayashi\endcsname{%
  \draw[lpcornerA] (0,13) -- (0,3) -- (3,0) -- (13,0); \draw[lpcornerB] (1.7,9.5) -- (1.7,3.7) -- (3.7,1.7) -- (9.5,1.7);
  \fill[Fac] (-0.6,13) rectangle (0.6,14.6); \fill[Fac] (13,-0.6) rectangle (14.6,0.6);
  \fill[FacDark] (3.4,3.4) circle (0.45);}
\def\lp@corner@grubstake{\draw[lpcornerA] (0,13) -- (0,0) -- (13,0); \draw[lpcornerB] (1.7,9.5) -- (1.7,1.7) -- (9.5,1.7);
  \foreach \p in {(0,13.9),(13.9,0),(0,0)}{\fill[Fac] \p circle (0.85); \fill[FacLight] \p circle (0.3);}
  \fill[FacDark] (4,4) circle (0.4);}
\def\lp@corner@synthexis{\draw[lpcornerA, rounded corners=4.5mm] (0,13) -- (0,0) -- (13,0);
  \draw[lpcornerB, rounded corners=3mm] (1.7,9.5) -- (1.7,1.7) -- (9.5,1.7);
  \fill[Fac] (0,13.6) circle (0.75); \fill[Fac] (13.6,0) circle (0.75);
  \draw[FacDark, line width=0.35pt] (4.2,4.2) circle (0.9); \fill[FacDark] (5.9,3.1) circle (0.35); \fill[FacDark] (3.1,5.9) circle (0.25);}
\def\lp@corner@player{\draw[lpcornerA] (0,13) -- (0,0) -- (13,0); \draw[lpcornerB] (1.7,9.5) -- (1.7,1.7) -- (9.5,1.7);
  \fill[Fac] (-0.9,-0.9) rectangle (2.3,2.3); \fill[FacOn] (0.7,0.7) circle (0.45);
  \fill[Fac] (-0.6,12.4) rectangle (0.6,14.2); \fill[Fac] (12.4,-0.6) rectangle (14.2,0.6);}
% A blank page carries the book's seal, small and faint, as a printer's device.
\newcommand\lp@blankseal{%
  \begin{tikzpicture}[x=1mm,y=1mm]
    \useasboundingbox (0,0) rectangle (148,210);
    \begin{scope}[shift={(74,108)}, transparency group, opacity=0.13, color=LPInk]\lpemblem{neutral}{17}{0.55}\end{scope}
  \end{tikzpicture}}
\AddToShipoutPictureBG{%
  \AtPageLowerLeft{\lp@ground}%
  \iflpdecor\iflpnotab\else\AtPageLowerLeft{\lp@watermark}\AtPageLowerLeft{\lp@corners}\fi\fi
  \iflpblank\AtPageLowerLeft{\lp@blankseal}\fi
  \AtPageLowerLeft{\lp@tab}\global\lpnotabfalse\global\lpblankfalse}
% Book matter is one sequence: the folio counts from the cover, as a printed booklet's do.
\renewcommand\frontmatter{}\renewcommand\mainmatter{\clearpage}\renewcommand\backmatter{}

\newcommand\lp@chaptertitle{}
\pagestyle{fancy}
\fancyhf{}
\renewcommand\headrulewidth{0pt}
\newcommand\lp@headicon{\tikz[x=1mm,y=1mm,baseline=-0.95mm]{\begin{scope}[color=FacDark]\lpemblem{\lp@fac}{2.9}{0.17}\end{scope}}}
\newcommand\lp@headmark[1]{{\color{FacDark}\lpkicker{#1}}}
% The folio sits in a plate of the faction's colour, its shape the faction's: a chamfered plate (Kessler-Hayashi),
% a round seal (Grubstake), a rounded cell (Synthexis), a square plate (yours and the book's).
\tikzset{lpfolio/.style={fill=Fac, text=FacOn, inner xsep=1.3mm, inner ysep=0.9mm, minimum height=4.6mm, minimum width=6.4mm},
  lpfolio@neutral/.style={rectangle},
  lpfolio@player/.style={rectangle},
  lpfolio@kessler-hayashi/.style={chamfered rectangle, chamfered rectangle xsep=1.1mm},
  lpfolio@grubstake/.style={circle, inner sep=0.5mm, minimum size=6.2mm, double=Fac, draw=FacOn, line width=0.35pt, double distance=0.5pt},
  lpfolio@synthexis/.style={rectangle, rounded corners=2.2mm}}
\newcommand\lp@folio{\tikz[baseline=(n.base)]{\node[lpfolio, lpfolio@\lp@fac/.try] (n)
  {\lpsemibold\fontsize{7.3}{8}\selectfont\addfontfeature{Numbers=Lining}\thepage};}}
\newcommand\lp@footrule{\textcolor{Fac}{\rule[0.62ex]{9mm}{0.5pt}}\hspace{0.6mm}\textcolor{Fac}{\rule[0.42ex]{1.4mm}{1.4mm}}}
\newcommand\lp@footruleR{\textcolor{Fac}{\rule[0.42ex]{1.4mm}{1.4mm}}\hspace{0.6mm}\textcolor{Fac}{\rule[0.62ex]{9mm}{0.5pt}}}
\fancyhead[LE]{\lpusefaction{\lpchapterfaction}\lp@headicon\hspace{1.8mm}\lp@headmark{\leftmark}}
\fancyhead[RO]{\lpusefaction{\lpchapterfaction}\lp@headmark{\rightmark}\hspace{1.8mm}\lp@headicon}
\fancyfoot[LE]{\lpusefaction{\lpchapterfaction}\lp@folio\hspace{1.6mm}\lp@footruleR}
\fancyfoot[RO]{\lpusefaction{\lpchapterfaction}\lp@footrule\hspace{1.6mm}\lp@folio}
\fancypagestyle{plain}{\fancyhf{}%
  \fancyfoot[LE]{\lpusefaction{\lpchapterfaction}\lp@folio\hspace{1.6mm}\lp@footruleR}%
  \fancyfoot[RO]{\lpusefaction{\lpchapterfaction}\lp@footrule\hspace{1.6mm}\lp@folio}}
\fancypagestyle{lpbare}{\fancyhf{}}

% ---------------------------------------------------------------- headings
% A chapter is never numbered (the manual is read by subject, not in order). Sections in Titillium.
\setcounter{secnumdepth}{-2}
\newcommand\lp@lozenge[1]{\tikz[x=1mm,y=1mm,baseline=-0.5*#1mm+0.5ex]\fill (0,#1) -- (#1,0) -- (0,-#1) -- (-#1,0) -- cycle;}
\newcommand\lp@sectitle[1]{\textcolor{Fac}{\rule[0.12ex]{2.2mm}{2.2mm}}\hspace{2.2mm}#1\nobreak\hspace{2.6mm}%
  {\color{FacMid}\leaders\hrule height 0.6ex depth -0.48ex\hfill}\nobreak\hspace{0.9mm}{\color{Fac}\lp@lozenge{0.9}}}
\titleformat{\section}{\sffamily\bfseries\fontsize{12.5}{15}\selectfont\color{FacDark}}{}{0pt}{\lp@sectitle}
\titlespacing*{\section}{0pt}{14pt plus 3pt}{5pt}
\titleformat{\subsection}{\lpsemibold\fontsize{10}{12}\selectfont\color{LPInk}}{}{0pt}{}
\titlespacing*{\subsection}{0pt}{10pt plus 2pt}{3pt}
\titleformat{\subsubsection}[runin]{\lpsemibold\fontsize{9}{11}\selectfont\color{FacDark}}{}{0pt}{}[.]
\titlespacing*{\subsubsection}{0pt}{6pt}{0.6em}
\renewcommand\sectionmark[1]{\markright{#1}}
% A heading keeps a few lines of its text with it (the manual's first full render left headings at the foot of a page).
\RequirePackage{needspace}
\preto\section{\Needspace{7\baselineskip}}
\preto\subsection{\Needspace{5\baselineskip}}

% \lpchapter{title}{label}{faction}{kicker}: a chapter that is not a faction's opener.
\newcommand\lpchapter[4]{%
  \clearpage\gdef\lpchapterfaction{#3}\lpusefaction{#3}%
  \markboth{#1}{#1}\thispagestyle{plain}%
  \phantomsection\addcontentsline{toc}{chapter}{\protect\lptocentry{#3}{#1}}\ifstrempty{#2}{}{\label{#2}}%
  \global\lpdecortrue
  \AddToShipoutPictureBG*{\AtPageUpperLeft{\color{lp@\lp@fac}\raisebox{-5mm}{\rule{\paperwidth}{5mm}}}}%
  \vspace*{3mm}%
  \ifstrempty{#4}{}{{\color{FacDark}\lpkicker{#4}}\par\vspace{2.2mm}}%
  \noindent\lp@banner{#1}\par\vspace{7mm}\@afterheading}
% The banner: a ribbon of the faction's colour (the book's slate for its own chapters), its tail folded behind at the
% left with a swallowtail cut, the title on it, and the faction's emblem in a medallion at its right end.
\definecolor{LPSlate}{HTML}{3B4048}
\newcommand\lp@banner[1]{%
  \ifdefstring{\lp@fac}{neutral}{\colorlet{BanMain}{LPSlate}\colorlet{BanDark}{LPSlate!60!black}\colorlet{BanOn}{white}}
    {\colorlet{BanMain}{Fac}\colorlet{BanDark}{FacDark}\colorlet{BanOn}{FacOn}}%
  \begin{tikzpicture}[x=1mm,y=1mm]
    \useasboundingbox (0,-3) rectangle (\lptextwidthmm,13);% the tail hangs into the margin
    \fill[BanDark] (-4,-2.6) -- (9,-2.6) -- (9,8.4) -- (-4,8.4) -- (-0.6,2.9) -- cycle;
    \fill[BanDark!70!black] (4,0) -- (4,-2.6) -- (9,-2.6) -- cycle;
    \fill[BanMain] (4,0) rectangle (\lptextwidthmm-9,11);
    \draw[BanOn, opacity=0.55, line width=0.35pt] (5.2,1.2) -- (\lptextwidthmm-9,1.2) (5.2,9.8) -- (\lptextwidthmm-9,9.8);
    \node[anchor=base west, inner sep=0pt, text=BanOn, font=\sffamily\bfseries\fontsize{19}{20}\selectfont\addfontfeature{LetterSpace=1.5}]
      at (9,3.3) {#1};
    \fill[BanDark] (\lptextwidthmm-9,5.5) circle (8.9);
    \fill[LPPaper] (\lptextwidthmm-9,5.5) circle (8.1);
    \draw[BanMain, line width=0.6pt] (\lptextwidthmm-9,5.5) circle (7.3);
    \begin{scope}[shift={(\lptextwidthmm-9,5.5)}, color=BanDark]\lpemblem{\lp@fac}{6.2}{0.26}\end{scope}
  \end{tikzpicture}}

% \lp@coverimage{w}{h}{path}: the image scaled to cover a w by h box (the caller clips it), whatever its proportions.
\newcommand\lp@coverimage[3]{%
  \setbox\z@\hbox{\includegraphics{#3}}%
  \pgfmathsetmacro\lp@imgr{\wd\z@/\ht\z@}\pgfmathsetmacro\lp@boxr{#1/#2}%
  \ifdim\lp@imgr pt<\lp@boxr pt\includegraphics[width=#1]{#3}\else\includegraphics[height=#2]{#3}\fi}
% \lpopener{title}{label}{faction}{hero path}{epigraph}{kicker}: a faction's chapter opens on its hero, an event.
% The vista runs to the trim on three sides, the faction's band under it carries the name. An opener is a recto, so
% the lore spread that follows it starts on the verso overleaf (lpstory below).
\newlength\lp@vistah \setlength\lp@vistah{96mm}
\newlength\lp@bandh \setlength\lp@bandh{19mm}
\newcommand\lpopener[6]{%
  \lp@cleartoodd\gdef\lpchapterfaction{#3}\lpusefaction{#3}%
  \markboth{#1}{#1}\thispagestyle{lpbare}\global\lpnotabtrue%
  \phantomsection\addcontentsline{toc}{chapter}{\protect\lptocentry{#3}{#1}}\ifstrempty{#2}{}{\label{#2}}%
  \AddToShipoutPictureBG*{\AtPageLowerLeft{\begin{tikzpicture}[x=1mm,y=1mm]
    \useasboundingbox (0,0) rectangle (148,210);
    \begin{scope}
      \clip (0,210) rectangle (148,114);
      \node[inner sep=0pt, anchor=center] at (74,162) {\lp@coverimage{152mm}{98mm}{#4}};
    \end{scope}
    \fill[lp@#3] (0,114) rectangle (148,95);
    \fill[LPNight, opacity=0.55] (0,114) rectangle (148,113.4);
    \ifodd\value{page}\def\lp@x{16}\else\def\lp@x{15}\fi
    \node[anchor=south west, inner sep=0pt, text=lp@#3@on, font=\sffamily\bfseries\fontsize{27}{28}\selectfont]
      at (\lp@x,99.6) {#1};
    \ifstrempty{#6}{}{\node[anchor=north west, inner sep=0pt, text=lp@#3@on, opacity=0.85,
      font=\lpsemibold\fontsize{6.6}{8}\selectfont\addfontfeature{LetterSpace=14}] at (\lp@x,111.4) {\MakeUppercase{#6}};}
    \fill[LPNight, opacity=0.3] (121.6,103.7) circle (13.3);
    \fill[lp@#3@dark] (121,104.5) circle (13);
    \fill[lp@#3] (121,104.5) circle (12.1);
    \draw[lp@#3@on, line width=0.5pt] (121,104.5) circle (11.2);
    \begin{scope}[shift={(121,104.5)}, color=lp@#3@on]\lpemblem{#3}{9.4}{0.34}\end{scope}
  \end{tikzpicture}}}%
  \global\lpdecortrue
  \vspace*{\dimexpr 210mm-95mm-18mm+5mm-\topskip\relax}%
  \ifstrempty{#5}{}{\lpepigraph{#5}}%
  \@afterheading}
\newcommand\lpepigraph[1]{%
  {\fontsize{12}{15}\selectfont\itshape\color{LPInk}\textcolor{FacDark}{\char"201C}#1\textcolor{FacDark}{\char"201D}\par}%
  \vspace{3mm}\lp@ornrule{30}\par\vspace{4mm}}
% \lp@ornrule{width, mm}: a rule broken by a lozenge between two beads.
\newcommand\lp@ornrule[1]{\tikz[x=1mm,y=1mm,baseline=-0.6mm]{%
  \draw[Fac, line width=0.7pt] (0,0) -- ({#1/2-4.2},0) ({#1/2+4.2},0) -- (#1,0);
  \fill[Fac] ({#1/2},1.5) -- ({#1/2+1.5},0) -- ({#1/2},-1.5) -- ({#1/2-1.5},0) -- cycle;
  \fill[FacDark] ({#1/2-2.9},0) circle (0.5) ({#1/2+2.9},0) circle (0.5);
  \fill[Fac] (0,0) circle (0.55) (#1,0) circle (0.55);}}

% ---------------------------------------------------------------- contents
\newcommand\lptocentry[2]{\lpusefaction{#1}\textcolor{Fac}{\rule[0.05ex]{2.4mm}{2.4mm}}\hspace{2.6mm}#2}
\renewcommand*\l@chapter[2]{%
  \addpenalty{-\@highpenalty}\vskip 3.2pt plus 1pt
  {\sffamily\fontsize{10.5}{13}\selectfont\color{LPInk}\noindent#1\nobreak
    \leaders\hbox{\hbox to 2.4mm{\hss\color{LPRule}.\hss}}\hfill\nobreak
    {\lpsemibold\addfontfeature{Numbers=Lining}#2}\par}}
\renewcommand*\l@section[2]{}
\renewcommand\tableofcontents{%
  \clearpage\thispagestyle{lpbare}\gdef\lpchapterfaction{neutral}%
  \vspace*{14mm}{\color{LPMuted}\lpkicker{Landing Party}}\par\vspace{2.2mm}
  \begingroup\lpusefaction{neutral}\noindent\lp@banner{Contents}\par\endgroup\vspace{8mm}
  \@starttoc{toc}\clearpage}

% ---------------------------------------------------------------- lists, tables, figures
\setlist{leftmargin=1.15em, itemsep=1.5pt, topsep=3pt, parsep=0pt}
\setlist[itemize,1]{label={\textcolor{Fac}{\rule[0.32ex]{1.5mm}{1.5mm}}}}
\setlist[itemize,2]{label={\textcolor{FacDark}{\rule[0.45ex]{1.4mm}{0.5pt}}}}
\setlist[enumerate,1]{label={\sffamily\bfseries\color{FacDark}\arabic*}, leftmargin=1.3em}
\AtBeginEnvironment{longtable}{\sffamily\fontsize{7.4}{9.4}\selectfont\addfontfeature{Numbers=Lining}\arrayrulecolor{LPRule}}
\AtBeginEnvironment{tabular}{\sffamily\fontsize{7.4}{9.4}\selectfont\addfontfeature{Numbers=Lining}}
\RequirePackage{caption}
\captionsetup{font={sf,scriptsize}, labelfont={color=LPMuted}, textfont={color=LPMuted}, justification=raggedright,
  singlelinecheck=false, skip=4pt, labelformat=empty, labelsep=none}% no "Figure 3:": the book never cites one by number
% An A5 page takes a big figure at its top or foot with a few lines of text: let it, rather than leave a page of floats.
\renewcommand\topfraction{0.86}\renewcommand\bottomfraction{0.75}\renewcommand\textfraction{0.1}
\renewcommand\floatpagefraction{0.8}\setcounter{topnumber}{2}\setcounter{bottomnumber}{1}\setcounter{totalnumber}{3}
\setlength\textfloatsep{12pt plus 3pt minus 2pt}\setlength\floatsep{10pt plus 2pt minus 2pt}
\renewcommand\texttt[1]{{\fontspec{TitilliumWeb-SemiBold}[Path=\lpextdir/fonts/, Extension=.ttf]\color{FacDark}#1}}

% ---------------------------------------------------------------- the unit card
% \begin{lpunit}{faction}{name}{tag line}{portrait path or empty} ... \end{lpunit}; inside, in order:
% \lprole, \lpstat (a bar), \lpstattext (no bar), \lpera, \lpunitlower{counters}{countered by}, then the lore.
\newlength\lp@portw \setlength\lp@portw{31mm}
\newlength\lp@barw \setlength\lp@barw{31mm}
% Two cards to a page: a card with a four-line blurb must stay under half the text height (the voice pass of
% 2026-11-05 made the blurbs four lines, and at the old spacing each card ran a few millimetres over and sat alone).
\newenvironment{lpunit}[4]{%
  \par\addvspace{7pt}\noindent\begin{minipage}{\linewidth}\lpusefaction{#1}%
  \begin{tikzpicture}[x=1mm,y=1mm]
    \useasboundingbox (0,-1.2) rectangle (\lptextwidthmm,7.2);
    \fill[Fac] (0,-1.2) rectangle (\lptextwidthmm,-0.6);
    \fill[Fac] (0,-0.6) rectangle (2.2,7.2);
    \node[anchor=base west, inner sep=0pt, font=\sffamily\bfseries\fontsize{12.5}{13}\selectfont\addfontfeature{LetterSpace=3}]
      at (4.2,1.4) {\MakeUppercase{#2}};
    \node[anchor=base east, inner sep=0pt, text=FacDark] at (\lptextwidthmm,1.4) {\lpkicker{#3}};
  \end{tikzpicture}\par\vspace{2mm}%
  \noindent\begin{minipage}[t]{\lp@portw}\vspace{0pt}\lp@portrait{#4}{#2}\end{minipage}\hfill
  \begin{minipage}[t]{\dimexpr\linewidth-\lp@portw-3.5mm\relax}\vspace{-0.8mm}\raggedright}
 {\par\end{minipage}\par\addvspace{7pt}}
\newcommand\lp@portrait[2]{%
  \begin{tikzpicture}[x=1mm,y=1mm]
    \ifstrempty{#1}{\fill[FacLight] (0,0) rectangle (31,31);
      \node[text width=26mm, align=center, font=\sffamily\scriptsize, text=LPMuted] at (15.5,15.5) {#2\\portrait to come};}
     {\node[inner sep=0pt, anchor=south west] at (0,0) {\includegraphics[width=31mm,height=31mm]{#1}};}
    \draw[LPInk, line width=0.5pt] (0,0) rectangle (31,31);
    \fill[Fac] (0,0) -- (5,0) -- (0,5) -- cycle;
  \end{tikzpicture}}
\newcommand\lprole[1]{{\itshape\fontsize{9}{11.4}\selectfont #1\par}\vspace{2.2mm}}
\newcommand\lp@statlabel[1]{\makebox[13.5mm][l]{\color{LPMuted}\lpsemibold\fontsize{6}{7}\selectfont\addfontfeature{LetterSpace=10}\MakeUppercase{#1}}}
\newcommand\lp@statvalue[1]{{\sffamily\fontsize{7.6}{9}\selectfont\addfontfeature{Numbers=Lining}#1}}
% \lpstat{label}{fraction of the faction's top, 0..1}{value}: the bar is cut in tenths, like a 90s readout.
\newcommand\lpstat[3]{\noindent\lp@statlabel{#1}%
  \begin{tikzpicture}[x=\lp@barw,y=1mm,baseline=0.1mm]
    \fill[FacLight] (0,0) rectangle (1,1.9);
    \fill[Fac] (0,0) rectangle (#2,1.9);
    \foreach \i in {0.1,0.2,...,0.95} \draw[white, line width=0.7pt] (\i,-0.1) -- (\i,2.0);
    \draw[FacDark, line width=0.35pt] (0,0) -- (1,0);
  \end{tikzpicture}\hspace{2mm}\lp@statvalue{#3}\par\vspace{0.35mm}}
\newcommand\lpstattext[2]{\noindent\lp@statlabel{#1}\lp@statvalue{#2}\par\vspace{0.35mm}}
% \lpera{label}{level or -1}{steps}{text}: pips for a corporation's eras, a gauge for an evolution (steps 1).
\newcommand\lpera[4]{\noindent\lp@statlabel{#1}%
  \ifdim #2pt<0pt\else
  \begin{tikzpicture}[x=1mm,y=1mm,baseline=0.1mm]
    \ifnum#3>1\relax
      \foreach \i in {0,...,#3}{\ifdim\i pt>#2pt\fill[FacLight] (\i*3.2,0) rectangle (\i*3.2+2.4,1.9);
        \else\fill[Fac] (\i*3.2,0) rectangle (\i*3.2+2.4,1.9);\fi}
    \else
      \fill[FacLight] (0,0) rectangle (\lp@barw,1.9);
      \fill[Fac] ({#2*\lp@barw/1mm},-0.5) rectangle ({#2*\lp@barw/1mm+0.9},2.4);
      \draw[FacDark, line width=0.35pt] (0,0) -- (\lp@barw,0);
    \fi
  \end{tikzpicture}\hspace{2mm}\fi\lp@statvalue{#4}\par\vspace{0.35mm}}
\newcommand\lpunitlower[2]{\par\end{minipage}\par\vspace{1.8mm}\noindent
  \begin{minipage}[t]{0.5\linewidth}\raggedright\lp@vs{Counters}{#1}\end{minipage}\hfill
  \begin{minipage}[t]{0.47\linewidth}\raggedright\lp@vs{Countered by}{#2}\end{minipage}\par\vspace{1.5mm}%
  {\color{LPRule}\rule{\linewidth}{0.4pt}}\par\vspace{1.2mm}%
  \begin{minipage}{\linewidth}\fontsize{9}{12}\selectfont}
\newcommand\lp@vs[2]{{\color{FacDark}\lpkicker{#1}}\\[0.6mm]{\fontsize{8.6}{10.8}\selectfont #2}}

% ---------------------------------------------------------------- trees
% The filter lays a tree out (lp-manual.lua, tree_tex) and draws it with these styles; nothing here places a node.
\tikzset{
  lpedge/.style={draw=LPInk!60, line width=0.55pt, rounded corners=1.2mm, -{Stealth[length=1.5mm,width=1.3mm]}},
  lpband/.style={fill=FacLight},
  lptlabel/.style={inner sep=0pt, align=center, font=\lpsemibold\fontsize{6.2}{7}\selectfont, text=LPInk},
  lptnote/.style={inner sep=0pt, align=center, font=\sffamily\fontsize{5.4}{6.2}\selectfont\addfontfeature{Numbers=Lining}, text=LPMuted},
  lptans/.style={inner sep=0pt, align=center, font=\sffamily\fontsize{5.4}{6.2}\selectfont, text=FacDark},
  lpbandlabel/.style={inner sep=0pt, rotate=90, font=\lpsemibold\fontsize{5.6}{6.4}\selectfont\addfontfeature{LetterSpace=12}, text=FacDark},
  lpcolhead/.style={inner sep=0pt, align=center, font=\lpsemibold\fontsize{6}{7}\selectfont\addfontfeature{LetterSpace=8}, text=LPInk},
}
\newenvironment{lptree}[3]{% faction, kicker, title
  \par\addvspace{8pt}\noindent\begin{minipage}{\linewidth}\lpusefaction{#1}%
  {\color{FacDark}\lpkicker{#2}}\par\vspace{0.6mm}
  {\sffamily\bfseries\fontsize{12.5}{15}\selectfont #3\par}\vspace{1.6mm}{\color{Fac}\rule{\linewidth}{0.8pt}}\par\vspace{2.4mm}%
  \noindent}
 {\par\end{minipage}\par\addvspace{8pt}}
% \lpbadge{colour}{text colour}{letter}: how a technology is got, a square stamp (no TikZ: it sits inside nodes).
\newcommand\lpbadge[3]{{\setlength\fboxsep{0.35mm}\colorbox{#1}{\makebox[1.5mm]{\textcolor{#2}{\lpsemibold\fontsize{4.8}{5}\selectfont #3}}}}}
\newcommand\lptreecaption[1]{\par\vspace{2mm}{\sffamily\fontsize{7}{9}\selectfont\color{LPMuted}#1\par}}

% ---------------------------------------------------------------- figures: annotated plans and sequence strips
% The filter places each marker from the plan's JSON (lp-manual.lua, callouts_latex); these are its looks.
% A marker: the faction's disc, ringed in white over a soft dark halo so it reads on sand, grass or creep.
\tikzset{
  lpmark/.style={circle, fill=Fac, draw=white, line width=0.75pt, inner sep=0pt, minimum size=4.1mm, text=FacOn,
    font=\lpsemibold\fontsize{6.6}{6.6}\selectfont\addfontfeature{Numbers=Lining}},
  lpmarkhalo/.style={circle, fill=black, opacity=0.38, inner sep=0pt, minimum size=5.3mm},
}
\newcommand\lpmark[3]{\node[lpmarkhalo] at (#1,#2) {}; \node[lpmark] at (#1,#2) {#3};}
% \lpleader{x}{y}{marker x}{marker y}: a nudged marker's line back to its point, and the point's dot.
\newcommand\lpleader[4]{\draw[black, opacity=0.45, line width=1.5pt] (#1,#2) -- (#3,#4);
  \draw[white, line width=0.6pt] (#1,#2) -- (#3,#4); \fill[white] (#1,#2) circle (0.75); \fill[Fac] (#1,#2) circle (0.45);}
% Plans and strips float to the top or foot of a page near where they are cited, as the stills do (fig-pos in
% _quarto.yml): a big picture held in place left half pages empty.
\newenvironment{lpcallouts}[1]{% faction
  \begin{figure}[!htb]\lpusefaction{#1}\noindent}
 {\par\end{figure}}
% A legend entry: three to a row, a long label hanging under itself.
\newcommand\lpcalloutitem[2]{%
  \parbox[t]{\dimexpr(\linewidth-4mm)/3\relax}{\raggedright\hangindent=4.9mm\hangafter=1\noindent
    \makebox[4.9mm][l]{\tikz[baseline=-0.75mm]\node[lpmark, minimum size=3.5mm, line width=0pt, draw=none,
      font=\lpsemibold\fontsize{6}{6}\selectfont\addfontfeature{Numbers=Lining}] {#1};}%
    {\sffamily\fontsize{7.2}{8.8}\selectfont\addfontfeature{Numbers=Lining}#2}\par\vspace{0.6mm}}\hskip 2mm\relax}
% \begin{lpstrip}{faction}{frames} \lpstripframe{n}{path}{caption}\hfill... \end{lpstrip}: frames side by side.
\newcount\lp@frames
\newenvironment{lpstrip}[2]{%
  \begin{figure}[!htb]\lpusefaction{#1}\global\lp@frames=#2\relax\noindent}
 {\par\end{figure}}
\newcommand\lpstripframe[3]{%
  \begin{minipage}[t]{\dimexpr(\linewidth-2.5mm*\numexpr\lp@frames-1\relax)/\lp@frames\relax}
    \vspace{0pt}\includegraphics[width=\linewidth]{#2}\par\vspace{1.3mm}
    {\color{Fac}\rule{\linewidth}{0.9pt}}\par\vspace{0.9mm}
    {\sffamily\fontsize{6.9}{8.6}\selectfont\raggedright\addfontfeature{Numbers=Lining}%
      {\lpsemibold\color{FacDark}#1}\hspace{1.2mm}#3\par}
  \end{minipage}}

% ---------------------------------------------------------------- lore
% A lore page is a found document and runs as a spread: it starts on a verso (a blank, bare page first if the next page
% is a recto), prints on the paper tint and ends its page, so the booklet opens flat on it. A dossier block (lpdossier)
% heads it: the Freight Authority's file on the faction, a field list set in the display face.
\newcommand\lp@cleartoeven{\clearpage\ifodd\value{page}\thispagestyle{lpbare}\global\lpnotabtrue\global\lpblanktrue\null\clearpage\fi}
\newcommand\lp@cleartoodd{\clearpage\ifodd\value{page}\else\thispagestyle{lpbare}\global\lpnotabtrue\global\lpblanktrue\null\clearpage\fi}
\newenvironment{lpstory}[3]{% kicker, title, draft (empty or a word)
  \lp@cleartoeven\global\lploretrue
  \vspace*{4mm}%
  % No kicker over a lore page's title by default (the user, 2026-11-06: "BACKGROUND" over "The story so far" read as
  % machine-made); the draft stamp stays.
  \ifstrempty{#1}{}{{\color{FacDark}\lpkicker{#1}}\par\vspace{1.4mm}}\ifstrempty{#3}{}{{\color{white}\colorbox{LPMuted}{\lpkicker{#3}}}\par\vspace{1.4mm}}
  {\fontsize{21}{25}\selectfont\itshape #2\par}\vspace{2.6mm}%
  \noindent\begin{tikzpicture}[x=1mm,y=1mm]
    \draw[LPInk!70, line width=0.5pt] (0,0) -- (52,0); \draw[LPInk!70, line width=0.5pt] (65,0) -- (\lptextwidthmm,0);
    \foreach \x in {55.3,58.5,61.7} \fill[Fac] (\x,0) +(0,1.1) -- +(1.1,0) -- +(0,-1.1) -- +(-1.1,0) -- cycle;
  \end{tikzpicture}\par\vspace{5mm}%
  \setlength{\parindent}{1.2em}\setlength{\parskip}{0pt}\fontsize{10}{14}\selectfont}
 {\par\clearpage\global\lplorefalse}
% The dossier: the file's fields (REGISTERED, PRINCIPALS, ...) as a two-column list on the faction's tint, in the
% display face, ruled top and bottom like a form. A label wider than its column pushes its text to the next line.
\newenvironment{lpdossier}{%
  \par\begin{tcolorbox}[enhanced, arc=0pt, outer arc=0pt, boxrule=0pt, colback=FacLight, left=3mm, right=3mm,
    top=2.8mm, bottom=2.4mm, before skip=1mm, after skip=9pt,
    overlay={\draw[FacDark, line width=0.9pt] (frame.north west) -- (frame.north east);
      \draw[FacDark, line width=0.4pt] (frame.south west) -- (frame.south east);
      \foreach \c in {north west, north east}\fill[FacDark] ([yshift=0.45pt]frame.\c) ++(0,-1.2mm) -- ++(1.2mm,1.2mm) -- ++(-2.4mm,0) -- cycle;
      \begin{scope}[x=1mm, y=1mm, shift={([xshift=-10mm,yshift=-0.5mm]frame.north east)}, rotate=-14,
          transparency group, opacity=0.2, color=FacDark]
        \lpemblem{\lp@fac}{7.6}{0.4}\draw[line width=0.8pt] (0,0) circle (8.8);
      \end{scope}}]
  \sffamily\fontsize{7.5}{9.4}\selectfont\addfontfeature{Numbers=Lining}\setlength{\parindent}{0pt}\setlength{\parskip}{0pt}%
  \begin{description}[leftmargin=27mm, labelwidth=25mm, labelsep=2mm, style=nextline, itemsep=2pt, parsep=0pt, topsep=0pt,
    font=\lpsemibold\fontsize{6.3}{8}\selectfont\addfontfeature{LetterSpace=10}\color{FacDark}]}
 {\end{description}\end{tcolorbox}\par}
% \lpdropcap{letter}{rest of the word}: the lore's first letter in a framed plate of the faction's colour, three
% lines deep, the rest of the word in small capitals (the illuminated capitals of the 90s manuals' history pages).
\newcommand\lpdropcap[2]{%
  \setbox\z@\hbox{\lp@capplate{#1}}%
  \hangindent=\dimexpr\wd\z@+2.4mm\relax\hangafter=-3
  \noindent\llap{\raisebox{-2\baselineskip}[0pt][0pt]{\box\z@}\hspace{2.4mm}}{\scshape #2}}
\newcommand\lp@capplate[1]{%
  \pgfmathsetmacro\lp@caph{(2*\the\baselineskip+\the\fontcharht\font`\H)/1mm}% two lines and a cap height
  \begin{tikzpicture}[x=1mm,y=1mm,baseline=0pt]
    \useasboundingbox (0,0) rectangle (\lp@caph,\lp@caph);
    \fill[FacDark] (0,0) rectangle (\lp@caph,\lp@caph);
    \fill[Fac] (0.55,0.55) rectangle (\lp@caph-0.55,\lp@caph-0.55);
    \draw[FacOn, line width=0.4pt] (1.3,1.3) rectangle (\lp@caph-1.3,\lp@caph-1.3);
    \foreach \x/\y in {1.3/1.3, 1.3/\lp@caph-1.3, \lp@caph-1.3/1.3, \lp@caph-1.3/\lp@caph-1.3}
      \fill[FacOn] (\x,\y) +(0,0.8) -- +(0.8,0) -- +(0,-0.8) -- +(-0.8,0) -- cycle;
    \begin{scope}[shift={(\lp@caph/2,\lp@caph/2)}, color=FacOn, opacity=0.16]\lpemblem{\lp@fac}{\lp@caph*0.42}{0.2}\end{scope}
    \node[text=FacOn, inner sep=0pt, font=\sffamily\bfseries\fontsize{30}{30}\selectfont] at (\lp@caph/2,\lp@caph/2-0.2) {#1};
  \end{tikzpicture}}
\setcounter{DefaultLines}{3}
\renewcommand\LettrineFontHook{\sffamily\bfseries\color{FacDark}}
\renewcommand\LettrineTextFont{\scshape}
\setlength\DefaultFindent{1.2pt}
\setlength\DefaultNindent{0pt}
% On a lore page it closes the history in less room (one rule, a smaller hand), so the page holds it and the chapter
% starts on the facing page (2026-11-06: the full quote alone on a recto read as a page left empty).
\newenvironment{lppullquote}[1]{% attribution or empty
  \par\begingroup\centering\def\lp@cite{#1}%
  \iflplore\vspace{4pt}\lp@ornrule{36}\par\nopagebreak\vspace{2mm}\nopagebreak\else\vspace{9pt}\lp@ornrule{36}\par\vspace{3.5mm}\fi
  \begin{minipage}{0.84\linewidth}\centering\setlength{\parindent}{0pt}%
  \iflplore\fontsize{11.5}{14.5}\else\fontsize{13.5}{17.5}\fi\selectfont\itshape\color{LPInk}}
 {\par\end{minipage}\par
  \ifdefempty{\lp@cite}{}{\nopagebreak\vspace{1.5mm}\nopagebreak{\color{FacDark}\lpkicker{\lp@cite}}\par}%
  \iflplore\else\vspace{3.5mm}\lp@ornrule{36}\par\vspace{9pt}\fi\endgroup\noindent\ignorespacesafterend}

% ---------------------------------------------------------------- sidebars
% The tip: a plate framed twice, riveted at the corners, its title on a tab with a cut corner standing on the frame.
\newtcolorbox{lptip}[1]{enhanced, arc=0pt, outer arc=0pt, boxrule=0.6pt,
  colback=FacLight, colframe=FacDark, left=3.6mm, right=3.4mm, top=3.6mm, bottom=2.8mm, before skip=21pt, after skip=9pt,
  fontupper=\sffamily\fontsize{8.4}{11}\selectfont\addfontfeature{Numbers=Lining},
  overlay={\draw[Fac, line width=0.35pt] ([shift={(0.9mm,0.9mm)}]frame.south west) rectangle ([shift={(-0.9mm,-0.9mm)}]frame.north east);
    \foreach \c/\dx/\dy in {south west/1/1, south east/-1/1, north west/1/-1, north east/-1/-1}
      \fill[FacDark] ([shift={(\dx*0.9mm,\dy*0.9mm)}]frame.\c) circle (0.55mm);
    \node[anchor=south west, chamfered rectangle, chamfered rectangle corners=north east, chamfered rectangle xsep=1.4mm,
      fill=Fac, text=FacOn, inner xsep=2mm, inner ysep=1.15mm, draw=FacDark, line width=0.6pt]
      at ([xshift=3.2mm, yshift=-0.3pt]frame.north west) {\lpkicker{\ifstrempty{#1}{Tip}{Tip: #1}}};},
  before upper={\setlength{\parskip}{3pt}}}
% The field note: a slip of paper under a paper clip, its edge dashed, its title on an ink tab with a swallowtail end.
% \begin{lpfieldnote}[prefix]{title}: the tab reads "prefix: title" (prefix "Field note"; an empty prefix, the title).
\newcommand\lp@notetab[2]{\ifstrempty{#1}{#2}{\ifstrempty{#2}{#1}{#1: #2}}}
\newtcolorbox{lpfieldnote}[2][Field note]{enhanced, arc=0pt, outer arc=0pt, boxrule=0pt,
  colback=LPPaper, left=4mm, right=4mm, top=5mm, bottom=3mm, before skip=13pt, after skip=8pt,
  fontupper=\fontsize{9.2}{12.4}\selectfont\itshape,
  overlay={\draw[LPInk!55, line width=0.4pt, dash pattern=on 1.6pt off 1.2pt] (frame.south west) rectangle (frame.north east);
    \node[anchor=west, fill=LPInk, text=white, inner xsep=1.8mm, inner ysep=1.1mm] (lp@tab) at ([xshift=3mm]frame.north west)
      {\lpkicker{\lp@notetab{#1}{#2}}};
    \fill[LPInk] (lp@tab.north east) -- ++(2.6mm,0) -- ([xshift=1.2mm]lp@tab.east) -- ([xshift=2.6mm]lp@tab.south east)
      -- (lp@tab.south east) -- cycle;
    \begin{scope}[x=1mm, y=1mm, shift={([xshift=-9mm, yshift=2.8mm]frame.north east)}, rotate=-8]
      \draw[LPInk!70, line width=0.75pt, rounded corners=1.1mm] (0.9,-4.8) -- (0.9,1.2) -- (3.3,1.2) -- (3.3,-7.4)
        -- (-0.3,-7.4) -- (-0.3,0.2) -- (2.1,0.2) -- (2.1,-3.6);
    \end{scope}},
  before upper={\setlength{\parskip}{3pt}}}

% ---------------------------------------------------------------- covers
% \lpcoverstrip{faction}{vista}{motto}: a band of the front cover. \lpcoverfactions is defined by the filter.
\providecommand\lpcoverfactions{}
% The front cover, as a box manual's (2026-11-06; the user: "The cover feels really underbaked compared to the rest of
% the manual"): the key art to the trim (\lpcoverart, a frame of the trailer), shaded top and bottom; a double brass
% frame with diamond corners; the title in two lines with a drop shadow; a ribbon for the subtitle; a crest for each
% faction (\lpcovercrests, its emblem in a medallion of its colours); the edition line at the foot.
\definecolor{LPBrass}{HTML}{C9A35A}
\definecolor{LPBrassDark}{HTML}{7E5F28}
\newcommand\lp@cx{0}
\newcommand\lp@covertitle[3]{% y, text, size
  \node[text=black, opacity=0.65, font=\sffamily\bfseries\fontsize{#3}{#3}\selectfont\addfontfeature{LetterSpace=7}]
    at (74.7,#1-0.8) {\MakeUppercase{#2}};
  \node[text=white, font=\sffamily\bfseries\fontsize{#3}{#3}\selectfont\addfontfeature{LetterSpace=7}]
    at (74,#1) {\MakeUppercase{#2}};}
\newcommand\lp@diamond[3]{\fill (#1,#2) +(0,#3) -- +(#3,0) -- +(0,-#3) -- +(-#3,0) -- cycle;}
\newcommand\lpfrontcover[3]{% title, subtitle, small line
  \thispagestyle{lpbare}\gdef\lpchapterfaction{neutral}%
  \AddToShipoutPictureBG*{\AtPageLowerLeft{\begin{tikzpicture}[x=1mm,y=1mm]
    \useasboundingbox (0,0) rectangle (148,210);
    \fill[LPNight] (-3,-3) rectangle (151,213);
    \ifdefempty{\lpcoverart}{}{\begin{scope}\clip (-3,-3) rectangle (151,213);
      \node[inner sep=0pt] at (74,105) {\lp@coverimage{156mm}{218mm}{\lpcoverart}};\end{scope}}
    \fill[LPNight, opacity=0.15] (-3,-3) rectangle (151,213);
    \fill[LPNight, opacity=0.94, path fading=south] (-3,118) rectangle (151,213);
    \fill[LPNight, opacity=0.97, path fading=north] (-3,-3) rectangle (151,104);
    \fill[LPNight, opacity=0.45] (-3,-3) rectangle (151,62);
    % the frame
    \draw[LPBrass, line width=0.9pt] (6,6) rectangle (142,204);
    \draw[LPBrass, line width=0.3pt] (7.7,7.7) rectangle (140.3,202.3);
    \foreach \x/\y in {6/6, 142/6, 6/204, 142/204} {\begin{scope}[color=LPBrass]\lp@diamond{\x}{\y}{2.6}\end{scope}
      \begin{scope}[color=LPNight]\lp@diamond{\x}{\y}{1.1}\end{scope}}
    % the title
    \lp@covertitle{184}{Landing}{44}
    \lp@covertitle{166.5}{Party}{44}
    % the ribbon: tails folded behind at both ends, swallowtail cuts, the subtitle on the band
    \fill[LPBrassDark] (25,140.6) -- (40,140.6) -- (40,149.4) -- (25,149.4) -- (29,145) -- cycle;
    \fill[LPBrassDark] (123,140.6) -- (108,140.6) -- (108,149.4) -- (123,149.4) -- (119,145) -- cycle;
    \fill[black!55!LPBrassDark] (40,140.6) -- (35,142.6) -- (40,142.6) -- cycle;
    \fill[black!55!LPBrassDark] (108,140.6) -- (113,142.6) -- (108,142.6) -- cycle;
    \fill[LPBrass] (35,142.6) rectangle (113,152.4);
    \draw[LPBrassDark, line width=0.35pt] (36.2,143.6) -- (111.8,143.6) (36.2,151.4) -- (111.8,151.4);
    \node[text=LPNight, font=\fontsize{13}{15}\selectfont\itshape] at (74,147.4) {#2};
    % the crests
    \gdef\lp@cx{\lpcrestfirst}\lpcovercrests
    % the foot
    \begin{scope}[color=LPBrass]\lp@diamond{74}{17.2}{0.9}\end{scope}
    \draw[LPBrass, line width=0.35pt] (48,17.2) -- (71.5,17.2) (76.5,17.2) -- (100,17.2);
    \node[text=LPBrass, inner sep=0pt] at (74,12.6) {\lpkicker{#3}};
  \end{tikzpicture}}}%
  \null\clearpage}
% \lpcovercrest{faction}: a faction's medallion on the front cover, its name under it; \lp@cx steps right.
\newcommand\lpcovercrest[1]{%
  \fill[black, opacity=0.5] (\lp@cx+0.6,42.4) circle (12);
  \fill[LPBrass] (\lp@cx,43) circle (12);
  \fill[lp@#1@dark] (\lp@cx,43) circle (11.1);
  \fill[lp@#1] (\lp@cx,43) circle (10.2);
  \draw[lp@#1@on, line width=0.5pt] (\lp@cx,43) circle (9.3);
  \begin{scope}[shift={(\lp@cx,43)}, color=lp@#1@on]\lpemblem{#1}{7.7}{0.3}\end{scope}
  \node[text=white, inner sep=0pt, font=\sffamily\bfseries\fontsize{6.4}{7}\selectfont\addfontfeature{LetterSpace=10}]
    at (\lp@cx,26.5) {\MakeUppercase{\csname lp@#1@name\endcsname}};
  \pgfmathsetmacro\lp@nx{\lp@cx+32}\xdef\lp@cx{\lp@nx}}
\providecommand\lpcrestfirst{42}
% \lpcoverstrip{faction}{vista}{motto}: one of the back cover's screenshots (the vistas, printed once), its faction's
% colour under it.
\newcommand\lpcoverstrip[3]{%
  \pgfmathsetmacro\lp@vw{(\lptextwidthmm-6)/3}%
  \begin{tikzpicture}[x=1mm,y=1mm]
    \begin{scope}\clip (0,0) rectangle (\lp@vw,27);
      \node[inner sep=0pt] at (\lp@vw/2,13.5) {\lp@coverimage{\lp@vw mm}{27mm}{#2}};\end{scope}
    \draw[white, opacity=0.35, line width=0.4pt] (0,0) rectangle (\lp@vw,27);
    \fill[lp@#1] (0,-1.6) rectangle (\lp@vw,0);
  \end{tikzpicture}}
\newcommand\lpinsidecover[1]{\thispagestyle{lpbare}\null\vfill{\sffamily\fontsize{7}{9.5}\selectfont\color{LPMuted}#1\par}\clearpage}
% The booklet is folded from sheets of four pages: blank pages before the back cover make the count a multiple
% of four, so the back cover is the last page of the last sheet.
\newcommand\lpbackcover[2]{% blurb, small print
  \clearpage\gdef\lpchapterfaction{neutral}\global\lpdecorfalse
  \lp@padtofour
  \thispagestyle{lpbare}%
  \AddToShipoutPictureBG*{\AtPageLowerLeft{\begin{tikzpicture}[x=1mm,y=1mm]
    \useasboundingbox (0,0) rectangle (148,210);
    \fill[LPNight] (-3,-3) rectangle (151,213);
  \end{tikzpicture}}}%
  \vspace*{6mm}{\color{white}\raggedright
    {\fontsize{22}{26}\selectfont\itshape Landing Party\par}\vspace{4mm}
    {\fontsize{11}{15.5}\selectfont #1\par}\vspace{7mm}
    \lpcoverfactions\vspace{6mm}
    \lpbackfactions
    \vfill{\sffamily\fontsize{6.6}{8.6}\selectfont\color{white!55}#2\par}}\null}
\providecommand\lpbackfactions{}
\RequirePackage{intcalc}
\newcommand\lp@padtofour{\ifnum\intcalcMod{\the\numexpr\ReadonlyShipoutCounter+1\relax}{4}=0
  \else\thispagestyle{lpbare}\global\lpblanktrue\null\clearpage\expandafter\lp@padtofour\fi}
% \lpbackfaction{faction}{portrait}{motto}: one of the four side by side on the back cover, the faction's own unit on
% the bestiary's grey ground (the cover vistas are printed once, on the front).
\newcommand\lpbackfaction[3]{%
  \begin{minipage}[t]{\dimexpr(\linewidth-9mm)/4\relax}
    \vspace{0pt}\includegraphics[width=\linewidth]{#2}\par\vspace{0pt}%
    {\color{lp@#1}\rule{\linewidth}{1.4mm}}\par\vspace{1.8mm}
    {\sffamily\bfseries\fontsize{7.4}{9}\selectfont\addfontfeature{LetterSpace=6}\MakeUppercase{\csname lp@#1@name\endcsname}\par}%
    \vspace{0.8mm}{\raggedright\itshape\fontsize{7.6}{9.4}\selectfont\color{white!72}#3\par}
  \end{minipage}}

\makeatother
